%% 1-2 10の まとまりを つくろう（10ページ、p.26〜35）
%% 第1巻のきまり：15 pt・行送り28 pt、分かち書き（文節の間に半角の空白）、漢字はすべてルビ
%% この章で定義した部品（頭文字 bkqb）：\bkqbstraw（ストロー）、\bkqbbundle（10本の たば）、\bkqbvpencil（たての えんぴつ）、
%%   \bkqbpbundle（えんぴつ 10本の たば）、\bkqbodo（メーターの かずの わ）、\bkqbhand（ひらいた て）、\bkqbbadge（A・B の まる）
\makeatletter
% ストロー  \bkqbstraw{(まん中)}{ながさ}{むき（度）}{色}
\newcommand\bkqbstraw[4]{\begin{scope}[shift={#1}, rotate=#3]
  \draw[#4, line width=1.9pt, line cap=round] (-#2/2,0) -- (#2/2,0);
  \draw[white, line width=0.45pt, opacity=0.7] (-#2/2+1,0.18) -- (#2/2-1,0.18);\end{scope}}
% ストロー 10本の たば  \bkqbbundle{(左下)}{ながさ}{色}   はば 約13 mm
\newcommand\bkqbbundle[3]{\begin{scope}[shift={#1}]
  \foreach \k in {0,...,9} {\draw[#3, line width=1.9pt, line cap=round] (0.9+\k*1.3,0) -- (0.9+\k*1.3,#2);
    \draw[white, line width=0.4pt, opacity=0.7] (1.1+\k*1.3,1) -- (1.1+\k*1.3,#2-1);}
  \fill[bkRed!85, rounded corners=0.5mm] (-0.4,#2/2-1.1) rectangle (13.1,#2/2+1.1);\end{scope}}
% たての えんぴつ  \bkqbvpencil{(下の まん中)}{ながさ}{色}   はば 1.6 mm
\newcommand\bkqbvpencil[3]{\begin{scope}[shift={#1}]
  \fill[#3] (-0.8,0) rectangle (0.8,#2-3.6);
  \fill[bkSkin] (-0.8,#2-3.6) -- (0.8,#2-3.6) -- (0.26,#2-0.9) -- (-0.26,#2-0.9) -- cycle;
  \fill[bkInk!85] (-0.28,#2-1.05) -- (0.28,#2-1.05) -- (0,#2) -- cycle;\end{scope}}
% えんぴつ 10本の たば  \bkqbpbundle{(左下)}{ながさ}{色}   はば 約18 mm
\newcommand\bkqbpbundle[3]{\begin{scope}[shift={#1}]
  \foreach \k in {0,...,9} {\bkqbvpencil{(0.9+\k*1.8,0)}{#2}{#3}}
  \fill[bkRed!85, rounded corners=0.5mm] (-0.3,#2/2-2.2) rectangle (18.3,#2/2-0.4);\end{scope}}
% メーターの かずの わ  \bkqbodo{(左下)}{すうじの ならび（例 1,2,3,9）}{この ばんごうから \R{右}{みぎ}の わを きいろに する（0なら ぬらない）}
\newcommand\bkqbodo[3]{\begin{scope}[shift={#1}]
  \fill[bkInk!85, rounded corners=1.5mm] (0,0) rectangle (38,13);
  \foreach \d [count=\k] in {#2} {
    \pgfmathtruncatemacro\bkqb@hl{(#3>0)&&(\k>=#3)}%
    \ifnum\bkqb@hl=1 \fill[bkYellow!60] (\k*8.6-6.2,2) rectangle ++(7.2,9);\else\fill[white] (\k*8.6-6.2,2) rectangle ++(7.2,9);\fi
    \node[font=\sffamily\bfseries\fontsize{16}{16}\selectfont, text=bkInk] at (\k*8.6-2.6,6.5) {\d};}\end{scope}}
% ひらいた て  \bkqbhand{(てのひらの まん中)}{1 または -1（-1で ひだりて）}
\newcommand\bkqbhand[2]{\begin{scope}[shift={#1}, xscale=#2]
  \foreach \a/\l in {18/7.4,6/8.4,-6/8.1,-18/7.1} {\filldraw[fill=bkSkin, draw=bkStem!70, line width=0.4pt, rounded corners=0.7mm, rotate=\a] (-0.8,0) rectangle (0.8,\l);}
  \filldraw[fill=bkSkin, draw=bkStem!70, line width=0.4pt, rounded corners=0.7mm, rotate around={55:(-2.2,-1.2)}] (-3.0,-1.2) rectangle (-1.4,4.0);
  \filldraw[fill=bkSkin, draw=bkStem!70, line width=0.4pt, rounded corners=1.5mm] (-3,-4.2) rectangle (3,2.6);\end{scope}}
% A・B の まる  \bkqbbadge{(x,y)}{文字}
\newcommand\bkqbbadge[2]{\node[circle, fill=bkYellow, text=white, font=\sffamily\bfseries\small, minimum size=5.5mm, inner sep=0pt] at #1 {#2};}
\makeatother

\bkchapter{1}{2}{10の まとまりを つくろう}{\bklevel{1}}{10の まとまりを つくろう}
\bkmichishirube{1-1}{1-3,1-4,1-5}

\begin{goal}
\begin{itemize}[label={\color{bkTheme}●}, leftmargin=1.4\zw, itemsep=0pt, topsep=0pt]
\item 10ずつ たばねて、 \R{大}{おお}きな かずを かぞえられる。
\item 23の 2と 3が、 なにを あらわすか わかる。
\item 100、 1000、 10000の かずが わかる。
\end{itemize}
\end{goal}

\begin{hatena}
\begin{bkfigure}
\begin{tikzpicture}[x=1mm, y=1mm]
  % A：ばらばら
  \fill[white, rounded corners=3mm] (0,0) rectangle (59,46);
  \bkqbbadge{(6,40)}{A}
  \foreach \x/\y/\a/\c in {21.3/13.9/117/bkRoute2,11.8/23.9/65/bkRoute4,25.1/31.5/22/bkRoute6,17.5/26.3/170/bkRoute3,30.9/20.3/175/bkRoute2,10.8/32.3/52/bkRoute4,14.5/13.1/55/bkRoute6,40.0/14.7/104/bkRoute3,39.2/28.2/43/bkRoute2,33.6/35.8/147/bkRoute4,19.8/20.0/120/bkRoute6,23.0/24.7/171/bkRoute3,35.2/23.4/111/bkRoute2,34.7/11.4/161/bkRoute4,45.1/23.7/26/bkRoute6,46.2/32.4/125/bkRoute3,10.1/17.3/46/bkRoute2,29.8/13.4/2/bkRoute4,34.4/30.4/161/bkRoute6,40.1/21.2/89/bkRoute3,19.0/35.0/175/bkRoute2,46.4/13.9/130/bkRoute4,28.8/27.7/169/bkRoute6} {
    \bkqbstraw{(\x+1.5,\y)}{19}{\a}{\c}}
  % B：10本ずつ たばねて ある
  \begin{scope}[shift={(63,0)}]
  \fill[white, rounded corners=3mm] (0,0) rectangle (59,46);
  \bkqbbadge{(6,40)}{B}
  \bkqbbundle{(8,12)}{19}{bkRoute6}
  \bkqbbundle{(25,12)}{19}{bkRoute6}
  \foreach \x in {44,48.5,53} {\bkqbstraw{(\x,21.5)}{19}{90}{bkRoute6}}
  \end{scope}
\end{tikzpicture}
\end{bkfigure}
\noindent A も B も、 ストローは おなじ かずだけ あります。\\
なん\R{本}{ぼん} あるか、 はやく わかるのは どっち？\\
どうして かな？
\end{hatena}

\clearpage
%% ---------------- p.27 やってみよう ----------------
\begin{yattemiyou}[10ずつ かこもう]
おはじきが たくさん。 なんこ あるかな？
\begin{enumerate}[label={\bkcircled[bkGreen]{\arabic*}}, leftmargin=2.2\zw, itemsep=2pt, topsep=2pt]
\item えんぴつで、 おはじきを 10こずつ ○で かこもう。 まちがえたら、 かこみなおせば いいよ。
\end{enumerate}
\begin{bkfigure}
\begin{tikzpicture}[x=1mm, y=1mm]
  \fill[white, rounded corners=3mm] (-3,-1) rectangle (123,48);
  \foreach \x/\y in {6.4/8.3,20.4/8.6,35.2/5.4,44.7/10.0,59.4/6.4,77.0/7.8,89.2/7.9,101.2/5.9,114.4/10.2,8.1/24.4,22.2/20.4,35.9/23.5,46.4/20.2,75.3/25.3,88.5/25.5,99.8/24.8,113.3/25.6,10.3/35.6,19.0/36.3,37.2/37.6,48.4/36.8,60.8/37.3,73.1/38.5,87.7/40.4,101.5/40.6,115.7/40.9} {
    \bkohajiki{(\x,\y)}{3}{bkRoute2}}
\end{tikzpicture}
\end{bkfigure}
\begin{enumerate}[label={\bkcircled[bkGreen]{\arabic*}}, leftmargin=2.2\zw, itemsep=4pt, topsep=0pt, start=2]
\item 10の まとまりは\enspace\bkkakikomi\enspace つ、\hspace{1\zw}ばらは\enspace\bkkakikomi\enspace こ
\item 「10、 20」と かぞえてから、 ばらを 「21、 22……」と かぞえよう。\par\vspace{2pt}
\noindent ぜんぶで\enspace\bkkakikomi\enspace こ
\item おうちでも、 ストローや ようじを 10\R{本}{ぽん}ずつ わゴムで たばねて みよう。
\end{enumerate}
\end{yattemiyou}

\bkfukidashi{sora}{10の まとまりを つくると、 かぞえるのが はやいし、 まちがえにくいね！}

\clearpage
%% ---------------- p.28〜29 数学の道具 ----------------
\bksection{doguu}{\R{数学}{すう|がく}の \R{道具}{どう|ぐ}}

\bksubsection{1}{10ずつ たばねる}

\label{bk:1-2:tabaneru}\bkhatenaref の B は、 10\R{本}{ぽん}の たばが 2つと、 ばらが 3\R{本}{ぼん}。

\begin{bkfigure}
\begin{tikzpicture}[x=1mm, y=1mm,
  cnt/.style={font=\sffamily\bfseries\small, text=bkTheme}]
  \bkqbbundle{(0,8)}{22}{bkRoute6}
  \bkqbbundle{(17,8)}{22}{bkRoute6}
  \foreach \x in {35,41.5,48} {\bkqbstraw{(\x,19)}{22}{90}{bkRoute6}}
  \node[cnt] at (6.5,3) {10};
  \node[cnt] at (23.5,3) {20};
  \foreach \x/\t in {35/21,41.5/22,48/23} {\node[cnt, font=\sffamily\bfseries\footnotesize] at (\x,3) {\t};}
  % 10の まとまりは どれも おなじ
  \begin{scope}[shift={(57,0)}]
  \fill[bkTheme!8, rounded corners=3mm] (0,-2) rectangle (66,33);
  \node[anchor=north, font=\sffamily\bfseries\small, text=bkTheme] at (33,32) {どれも 「10の まとまり」};
  \bkqbbundle{(4,4)}{18}{bkRoute6}
  \bktenframe{(21,8)}{4.2}{10}{bkRoute2}{0}{bkRoute2}
  \bkrod{(53,2)}{2.1}{bkRoute4}
  \node[font=\sffamily\footnotesize] at (10.5,0.6) {たば};
  \node[font=\sffamily\footnotesize] at (31.5,4.6) {10の わく};
  \node[font=\sffamily\footnotesize] at (60.5,8) {ぼう};
  \end{scope}
\end{tikzpicture}
\end{bkfigure}

\noindent 「10、 20」と かぞえて、 のこりは 「21、 22、 23」。 23\R{本}{ぼん}です。\\
だから、 \bkhatenaref の こたえは B。 たばねると、 すぐ わかるね。\bkidx{じゅうのまとまり}{10の まとまり}

\noindent どうして はやいのかな？ ばらばらの A は、 「1、 2、 3、 …、 23」と 23かい かぞえます。 とちゅうで どこまで かぞえたか わからなく なる ことも あります。 B は、 たばを 「10、 20」と 2かい、 ばらを 「21、 22、 23」と 3かい。 5かい いうだけで おわります。

\clearpage
\bksubsection{2}{10が いくつ？ なんじゅう？}

\label{bk:1-2:nanjuu}10の たばが 2つで 20（にじゅう）。 「に」は 2つ、 「じゅう」は 10の ことです。 たばを 1つずつ ふやして みよう。

\begin{bkfigure}
\begin{tikzpicture}[x=1mm, y=1mm]
  \foreach \n/\r [count=\k from 1] in {10/じゅう,20/にじゅう,30/さんじゅう,40/よんじゅう,50/ごじゅう,60/ろくじゅう,70/ななじゅう,80/はちじゅう,90/きゅうじゅう,100/ひゃく} {
    \pgfmathtruncatemacro\cx{mod(\k-1,5)}\pgfmathtruncatemacro\cy{div(\k-1,5)}
    \begin{scope}[shift={(\cx*26.4,-\cy*36)}]
      \fill[white, draw=bkLine, line width=0.8pt, rounded corners=2mm] (0,0) rectangle (24.4,34);
      \pgfmathsetmacro\st{12.2-(\k*2.2-0.6)/2}
      \foreach \j in {1,...,\k} {\bkrod{(\st+\j*2.2-2.2,14)}{1.6}{bkRoute4}}
      \node[font=\sffamily\bfseries\small, text=bkInk] at (12.2,9) {\n};
      \node[font=\fontsize{9}{10}\selectfont, text=bkInk] at (12.2,3.4) {\r};
    \end{scope}}
\end{tikzpicture}
\end{bkfigure}

\noindent 10が 3つで さんじゅう、 4つで よんじゅう…と、 いくつ あるかを 「じゅう」の まえに いいます。 10が 10こ そろうと、 100（ひゃく）に なります（\pageref{bk:1-2:hyaku}ページ）。

\noindent 10と ばらを あわせた かずも、 おなじように よみます。

\begin{bkfigure}
\begin{tikzpicture}[x=1mm, y=1mm]
  \foreach \b/\n/\r [count=\k from 0] in {1/11/じゅういち,3/13/じゅうさん,7/17/じゅうなな} {
    \begin{scope}[shift={(\k*44,0)}]
      \fill[white, draw=bkLine, line width=0.8pt, rounded corners=2mm] (0,0) rectangle (41,27);
      \bktenframe{(2.5,13)}{4.2}{10}{bkRoute2}{0}{bkRoute2}
      \pgfmathtruncatemacro\nb{\b}
      \foreach \j in {1,...,\nb} {\pgfmathsetmacro\bx{27+mod(\j-1,4)*3.6}\pgfmathsetmacro\by{19.3-div(\j-1,4)*3.6}\bkohajiki{(\bx,\by)}{1.5}{bkRoute4}}
      \node[font=\sffamily\bfseries\small] at (20.5,8) {10と \b で \n};
      \node[font=\fontsize{9}{10}\selectfont] at (20.5,3) {\r};
    \end{scope}}
\end{tikzpicture}
\end{bkfigure}

\noindent 10と 1で 11（じゅういち）、 10と 2で 12（じゅうに）…、 10と 9で 19（じゅうく）。 そして、 20と 3で 23（にじゅうさん）です。 \R{前}{まえ}の ほうで 10が いくつ、 あとの ほうで ばらが いくつかを いって いるね。

\bksubsection{3}{\R{十}{じゅう}の くらい、 \R{一}{いち}の くらい}

\label{bk:1-2:kurai}23の 2と 3は、 それぞれ なにを あらわして いるのかな？

\begin{bkfigure}
\begin{tikzpicture}[x=1mm, y=1mm]
  \fill[bkBlue!12, rounded corners=2mm] (0,0) rectangle (44,40);
  \fill[bkOrange!14, rounded corners=2mm] (46,0) rectangle (90,40);
  \node[font=\sffamily\bfseries\small, text=bkBlue!80!black] at (22,36) {\R{十}{じゅう}の くらい};
  \node[font=\sffamily\bfseries\small, text=bkOrange!85!black] at (68,36) {\R{一}{いち}の くらい};
  \bkrod{(15,10)}{2.2}{bkRoute4}\bkrod{(24,10)}{2.2}{bkRoute4}
  \foreach \x in {58,64,70} {\bkcube{(\x,10)}{2.6}{bkRoute4}}
  \node[font=\sffamily\bfseries\fontsize{24}{24}\selectfont, text=bkBlue!80!black] at (22,4.8) {2};
  \node[font=\sffamily\bfseries\fontsize{24}{24}\selectfont, text=bkOrange!85!black] at (68,4.8) {3};
  \node[anchor=west, font=\sffamily\bfseries\small, align=left] at (94,26) {2は 10が 2つ};
  \node[anchor=west, font=\sffamily\bfseries\small, align=left] at (94,14) {3は 1が 3つ};
\end{tikzpicture}
\end{bkfigure}

\noindent すうじを かく ばしょを 「くらい」と いいます。 おなじ 2でも、 \R{十}{じゅう}の くらいに あると 10が 2つ の ことです。\bkidx{くらい}{くらい}\bkidx{じゅうのくらい}{\R{十}{じゅう}の くらい}\bkidx{いちのくらい}{\R{一}{いち}の くらい}
10は、 10が 1つで ばらは なし。 ばらが ない ことを あらわす 「0」の ことは、 つぎの 1-3で くわしく まなびます。

\noindent お\R{金}{かね}でも おなじです。 10\R{円}{えん}\R{玉}{だま} 2まいと 1\R{円}{えん}\R{玉}{だま} 3まいで 23\R{円}{えん}。 10\R{円}{えん}\R{玉}{だま}の まいすうを \R{十}{じゅう}の くらいに、 1\R{円}{えん}\R{玉}{だま}の まいすうを \R{一}{いち}の くらいに かきます。

\begin{bkfigure}
\begin{tikzpicture}[x=1mm, y=1mm]
  \fill[bkBlue!12, rounded corners=2mm] (0,0) rectangle (44,30);
  \fill[bkOrange!14, rounded corners=2mm] (46,0) rectangle (90,30);
  \node[font=\sffamily\bfseries\small, text=bkBlue!80!black] at (22,26) {\R{十}{じゅう}の くらい};
  \node[font=\sffamily\bfseries\small, text=bkOrange!85!black] at (68,26) {\R{一}{いち}の くらい};
  \foreach \x in {16,28} {\bkcoin{(\x,14)}{10}}
  \foreach \x in {58,68,78} {\bkcoin{(\x,14)}{1}}
  \node[font=\sffamily\bfseries\fontsize{20}{20}\selectfont, text=bkBlue!80!black] at (22,4.5) {2};
  \node[font=\sffamily\bfseries\fontsize{20}{20}\selectfont, text=bkOrange!85!black] at (68,4.5) {3};
  \node[anchor=west, font=\sffamily\bfseries\small, align=left] at (94,20) {10\R{円}{えん}が 2まい};
  \node[anchor=west, font=\sffamily\bfseries\small, align=left] at (94,9) {1\R{円}{えん}が 3まい};
\end{tikzpicture}
\end{bkfigure}

\noindent もし くらいを まちがえて 32と かくと、 10\R{円}{えん}\R{玉}{だま}が 3まい、 1\R{円}{えん}\R{玉}{だま}が 2まいの ことに なって しまいます。 おなじ 2と 3でも、 かく へやで いみが かわるのです。

\clearpage
\bksubsection{4}{かずの ひょう（1から 100まで）}

\label{bk:1-2:hyou}1から 100までを、 10ずつ 1\R{行}{ぎょう}に ならべました。

\begin{bkfigure}
\begin{tikzpicture}[x=1mm, y=1mm]
  \foreach \r in {0,...,9} {\foreach \c in {1,...,10} {
    \pgfmathtruncatemacro\v{\r*10+\c}
    \ifnum\c=10 \def\fc{bkBlue!14}\else\def\fc{white}\fi
    \filldraw[fill=\fc, draw=bkGray!60, line width=0.5pt] ({(\c-1)*11.6},{-\r*7.6}) rectangle ++(11.6,7.6);
    \node[font=\sffamily\bfseries\fontsize{10}{10}\selectfont, text=bkInk] at ({(\c-1)*11.6+5.8},{-\r*7.6+3.8}) {\v};}}
  \draw[bkTheme, line width=1.4pt, rounded corners=1mm] (92.8,-22.8) rectangle (116,-15.2);
  \draw[bkGreen!80!black, line width=1.4pt, rounded corners=1mm] (92.8,-68.4) rectangle (116,-60.8);
  \node[anchor=west, font=\sffamily\bfseries\footnotesize, text=bkTheme, align=left] at (117,-19) {39の\\つぎ};
  \node[anchor=west, font=\sffamily\bfseries\footnotesize, text=bkGreen!70!black, align=left] at (117,-64.6) {99の\\つぎ};
\end{tikzpicture}
\end{bkfigure}

\noindent \R{右}{みぎ}へ 1つ すすむと 1 ふえます。 \R{下}{した}へ 1つ すすむと、 10の まとまりが 1つ ふえるので 10 ふえます。 いちばん \R{右}{みぎ}の れつは、 10、 20、 30…と、 ばらが ない かずです。

\noindent 39の つぎを みて みよう。 39は 10が 3つと ばらが 9こ。 ばらが 1こ ふえると 10こ そろうので、 たばねて 10の まとまりに します。 10が 3つと 1つで 4つ、 ばらは なし。 だから 39の つぎは 40です。 ひょうの あかい わくで たしかめよう。 みどりの わくの 99の つぎは、 \pageref{bk:1-2:hyaku}ページで しらべます。

\clearpage
\bksubsection{5}{100と \R{百}{ひゃく}の くらい}

\label{bk:1-2:hyaku}10の まとまりが 10こ あつまると、 100（ひゃく）。

\begin{bkfigure}
\begin{tikzpicture}[x=1mm, y=1mm]
  \foreach \k in {0,...,9} {\bkrod{(\k*2.2,6)}{1.6}{bkRoute4}}
  \draw[bk arrow, bkTheme] (24,14) -- (31,14);
  \bkflat{(33,6)}{1.6}{bkRoute4}
  \node[font=\sffamily\bfseries\footnotesize] at (10.5,2.6) {ぼう 10\R{本}{ぽん}};
  \node[font=\sffamily\bfseries\footnotesize] at (41,2.6) {100の いた};
  % 3つの へや
  \begin{scope}[shift={(54,0)}]
  \foreach \px/\c/\t in {0/bkGreen/{\R{百}{ひゃく}の くらい},23/bkBlue/{\R{十}{じゅう}の くらい},46/bkOrange/{\R{一}{いち}の くらい}} {
    \fill[\c!12, rounded corners=2mm] (\px,0) rectangle ++(22,32);
    \node[font=\sffamily\bfseries\fontsize{9}{10}\selectfont, text=\c!80!black] at (\px+11,29) {\t};}
  \bkflat{(1.8,8)}{1.45}{bkRoute4}\bkflat{(4.6,10.3)}{1.45}{bkRoute4}
  \foreach \x in {29,33,37} {\bkrod{(\x,8)}{1.6}{bkRoute4}}
  \foreach \x/\y in {51/9,56/9,61/9,66/9} {\bkcube{(\x,\y)}{2.2}{bkRoute4}}
  \foreach \px/\d/\c in {11/2/bkGreen,34/3/bkBlue,57/4/bkOrange} {\node[font=\sffamily\bfseries\fontsize{20}{20}\selectfont, text=\c!80!black] at (\px,3.6) {\d};}
  \end{scope}
\end{tikzpicture}
\end{bkfigure}

\noindent 234は、 100が 2つ、 10が 3つ、 1が 4つ。 100が いくつ あるかを かく へやを、 \R{百}{ひゃく}の くらいと いいます。 よみかたは 「にひゃく さんじゅう よん」です。\bkidx{ひゃくのくらい}{\R{百}{ひゃく}の くらい}

\noindent \pageref{bk:1-2:hyou}ページの ひょうの 99の つぎも、 これで わかります。 99は 10が 9つと ばらが 9こ。 ばらが 1こ ふえると 10こ そろうので、 たばねて 10が 10こ。 10が 10こ そろったので、 また まとめて 100の いた 1まい。 だから 99の つぎは 100です。 \R{百}{ひゃく}の くらいに 1、 \R{十}{じゅう}と \R{一}{いち}の くらいは なにも ないので 0を かいて、 すうじを 3つ ならべます。

\bksubsection{4}{1000と 10000}

\label{bk:1-2:ookii}10\R{円}{えん}\R{玉}{だま} 10まいで、 100\R{円}{えん}\R{玉}{だま} 1まいと おなじ。 お\R{金}{かね}にも 10の まとまりが あるよ。

\begin{bkfigure}
\begin{tikzpicture}[x=1mm, y=1mm,
  ar/.style={bk arrow, bkTheme, line width=1.1pt},
  lb/.style={font=\sffamily\bfseries\fontsize{8}{9}\selectfont, text=bkTheme},
  nm/.style={font=\sffamily\bfseries\small, text=bkInk},
  rd/.style={font=\fontsize{9}{10}\selectfont, text=bkInk}]
  \bkcoin{(4,10)}{1}
  \bkcoin{(30,10)}{10}
  \bkcoin{(56,10)}{100}
  \bkbill{(84,10)}{1000}
  \bkbill{(114,10)}{10000}
  \foreach \a/\b in {9/25,35/51,61/75,93/105} {\draw[ar] (\a,10) -- (\b,10); \node[lb] at ({(\a+\b)/2},14) {10まいで};}
  \foreach \x/\n/\r in {4/1/いち,30/10/じゅう,56/100/ひゃく,84/1000/せん,114/10000/いちまん} {
    \node[nm] at (\x,1) {\n}; \node[rd] at (\x,-4.6) {\r};}
\end{tikzpicture}
\end{bkfigure}

\noindent 100が 10こで 1000（せん）、 1000が 10こで 10000（いちまん）。 10こ あつまる たびに、 1つ \R{左}{ひだり}の くらいへ すすみます。\bkidx{せん}{1000（せん）}

\noindent 1000が いくつ あるかを かく へやを、 \R{千}{せん}の くらいと いいます。 \R{百}{ひゃく}の くらいの 1つ \R{左}{ひだり}です。\bkidx{せんのくらい}{\R{千}{せん}の くらい}

\begin{bkfigure}
\begin{tikzpicture}[x=1mm, y=1mm]
  \foreach \px/\c/\t/\d in {0/bkPurple/{\R{千}{せん}の くらい}/1,31/bkGreen/{\R{百}{ひゃく}の くらい}/2,62/bkBlue/{\R{十}{じゅう}の くらい}/4,93/bkOrange/{\R{一}{いち}の くらい}/3} {
    \fill[\c!12, rounded corners=2mm] (\px,0) rectangle ++(29,30);
    \node[font=\sffamily\bfseries\fontsize{9}{10}\selectfont, text=\c!80!black] at (\px+14.5,26.5) {\t};
    \node[font=\sffamily\bfseries\fontsize{20}{20}\selectfont, text=\c!80!black] at (\px+14.5,4) {\d};}
  \bkbill{(14.5,15)}{1000}
  \foreach \x in {40,52} {\bkcoin{(\x,15)}{100}}%
  \foreach \x/\y in {71.5/19,84.5/19,71.5/11,84.5/11} {\bkcoin{(\x,\y)}{10}}
  \foreach \x in {100,107.5,115} {\bkcoin{(\x,15)}{1}}
\end{tikzpicture}
\end{bkfigure}

\noindent 1000が 1つ、 100が 2つ、 10が 4つ、 1が 3つで 1243。 よみかたは 「せん にひゃく よんじゅう さん」です。

\begin{taisetsu}
10 あつまったら たばねて、 1つ \R{左}{ひだり}の くらいへ。\\
23は、 10が 2つと 1が 3つ。\\
くらいは \R{右}{みぎ}から、 \R{一}{いち}・\R{十}{じゅう}・\R{百}{ひゃく}・\R{千}{せん}。
\end{taisetsu}

%% （第11回：ここで ページを かえない）
%% ---------------- p.30〜31 例題 ----------------
\begin{reidai}[たばで かぞえる]
えんぴつは なん\R{本}{ぼん}？\par\vspace{3pt}
\noindent\begin{minipage}[c]{0.4\linewidth}
\kangaekata 10、 20、 30。 ばらが 5\R{本}{ほん}で 35。

\kotae 35\R{本}{ほん}
\end{minipage}\hfill
\begin{minipage}[c]{0.58\linewidth}\centering
\begin{tikzpicture}[x=1mm, y=1mm]
  \foreach \x in {0,20.5,41} {\bkqbpbundle{(\x,0)}{20}{bkRoute6}}
  \foreach \x in {0,1,2,3,4} {\bkqbvpencil{(63+\x*2.5,0)}{20}{bkRoute6}}
\end{tikzpicture}
\end{minipage}
\end{reidai}

\begin{jibun}
\noindent\begin{minipage}[c]{0.46\linewidth}
10\R{円}{えん}\R{玉}{だま}が 4まい、 1\R{円}{えん}\R{玉}{だま}が 7まい。 ぜんぶで なん\R{円}{えん}？
\end{minipage}\hfill
\begin{minipage}[c]{0.52\linewidth}\centering
\begin{tikzpicture}[x=1mm, y=1mm]
  \foreach \x in {0,8.5,17,25.5} {\bkcoin{(\x,0)}{10}}
  \foreach \x in {35,42,49,56} {\bkcoin{(\x,3.6)}{1}}
  \foreach \x in {38.5,45.5,52.5} {\bkcoin{(\x,-3.6)}{1}}
\end{tikzpicture}
\end{minipage}
\bkkotaeran[6\zw]
\bkjans{47\R{円}{えん}}
\end{jibun}

\begin{reidai}[くらいの へや]
56は、 10が いくつと 1が いくつ？\par\vspace{3pt}
\noindent\begin{minipage}[c]{0.56\linewidth}
\kangaekata \R{十}{じゅう}の くらいが 5、 \R{一}{いち}の くらいが 6。

\kotae 10が 5つと 1が 6つ
\end{minipage}\hfill
\begin{minipage}[c]{0.42\linewidth}\centering
\begin{tikzpicture}[x=1mm, y=1mm]
  \fill[bkBlue!12, rounded corners=2mm] (0,0) rectangle (25,27);
  \fill[bkOrange!14, rounded corners=2mm] (26,0) rectangle (51,27);
  \foreach \k in {0,...,4} {\bkrod{(2.5+\k*4.2,7)}{1.7}{bkRoute4}}
  \foreach \x/\y in {30/7,35/7,40/7,30/12,35/12,40/12} {\bkcube{(\x,\y)}{2.4}{bkRoute4}}
  \node[font=\sffamily\bfseries\large, text=bkBlue!80!black] at (12.5,3) {5};
  \node[font=\sffamily\bfseries\large, text=bkOrange!85!black] at (38.5,3) {6};
\end{tikzpicture}
\end{minipage}
\end{reidai}

\begin{jibun}
84は、 10が\enspace\bkkakikomi\enspace つと、 1が\enspace\bkkakikomi\enspace つ。
\bkjans{10が 8つと 1が 4つ}
\end{jibun}

\begin{reidai}[かずの ならび]
あいて いる ますに はいる かずは？\par\vspace{3pt}
\noindent\begin{minipage}[c]{0.46\linewidth}
\kangaekata 39の つぎは、 ばらが 10こ。 たばねて 10が 4つ。

\kotae 40
\end{minipage}\hfill
\begin{minipage}[c]{0.52\linewidth}\centering
\begin{tikzpicture}[x=1mm, y=1mm, nb/.style={draw=bkBlue, line width=0.8pt, fill=white, rounded corners=1.5mm, minimum width=11mm, minimum height=9mm, font=\sffamily\bfseries}]
  \foreach \x/\t in {0/37,13/38,26/39,52/41} {\node[nb] at (\x,0) {\t};}
  \node[nb, draw=bkTheme, line width=1.2pt] at (39,0) {40};
\end{tikzpicture}
\end{minipage}
\end{reidai}

\begin{jibun}
かずを かこう。\hspace{2\zw}68、\enspace 69、\enspace\bkkakikomi、\enspace 71
\bkjans{70}
\end{jibun}

\begin{reidai}[100より \R{大}{おお}きい かず]
100\R{円}{えん}\R{玉}{だま} 2まい、 10\R{円}{えん}\R{玉}{だま} 3まい、 1\R{円}{えん}\R{玉}{だま} 4まい。 なん\R{円}{えん}？\par\vspace{3pt}
\noindent\begin{minipage}[c]{0.5\linewidth}
\kangaekata 100が 2つ、 10が 3つ、 1が 4つ。

\kotae 234\R{円}{えん}
\end{minipage}\hfill
\begin{minipage}[c]{0.48\linewidth}\centering
\begin{tikzpicture}[x=1mm, y=1mm]
  \foreach \x in {0,9} {\bkcoin{(\x,0)}{100}}
  \foreach \x in {21,29.5,38} {\bkcoin{(\x,0)}{10}}
  \foreach \x in {48,55} {\bkcoin{(\x,3.6)}{1}\bkcoin{(\x,-3.6)}{1}}
\end{tikzpicture}
\end{minipage}
\end{reidai}

\begin{jibun}
100\R{円}{えん}\R{玉}{だま} 4まい、 10\R{円}{えん}\R{玉}{だま} 1まい、 1\R{円}{えん}\R{玉}{だま} 6まい。 なん\R{円}{えん}？\par\vspace{3pt}
\noindent\begin{minipage}[c]{0.6\linewidth}\centering
\begin{tikzpicture}[x=1mm, y=1mm]
  \foreach \x in {0,8.5,17,25.5} {\bkcoin{(\x,0)}{100}}
  \bkcoin{(36,0)}{10}
  \foreach \x in {46,53,60} {\bkcoin{(\x,3.6)}{1}\bkcoin{(\x,-3.6)}{1}}
\end{tikzpicture}
\end{minipage}\hfill
\begin{minipage}[c]{0.38\linewidth}
\bkkotaeran[6\zw]
\end{minipage}
\bkjans{416\R{円}{えん}}
\end{jibun}

%% ---------------- 第11回：千の くらい・たばが そろう ときの 例題 ----------------
\begin{reidai}[たばが 10こ そろう とき]
99の つぎの かずは？\par\vspace{3pt}
\noindent\begin{minipage}[c]{0.5\linewidth}
\kangaekata ばらが 9こと 1こで 10こ。 たばねると、 10の たばが 9つと 1つで 10こ。 10の たばが 10こで 100。

\kotae 100
\end{minipage}\hfill
\begin{minipage}[c]{0.48\linewidth}\centering
\begin{tikzpicture}[x=1mm, y=1mm]
  \foreach \k in {0,...,8} {\bkrod{(\k*2.3,0)}{1.6}{bkRoute4}}
  \foreach \cx/\cy in {0/0,1/0,2/0,0/1,1/1,2/1,0/2,1/2,2/2} {\bkcube{(23+\cx*2.2,\cy*2.2)}{1.6}{bkRoute4}}
  \bkcube{(30.5,0)}{1.6}{bkAppleRed}
  \draw[bk arrow, bkTheme] (34,8) -- (38,8);
  \bkflat{(40,0)}{1.6}{bkRoute4}
\end{tikzpicture}
\end{minipage}
\end{reidai}

\begin{jibun}
あいて いる ますに はいる かずを かこう。\par\vspace{2pt}
\noindent\hspace*{1\zw}（1）79、\enspace\bkkakikomi\hspace{4\zw}（2）109、\enspace\bkkakikomi
\bkjans{（1）80\quad（2）110}
\end{jibun}

\begin{reidai}[1000より \R{大}{おお}きい かず]
1000\R{円}{えん}さつ 1まい、 100\R{円}{えん}\R{玉}{だま} 2まい、 10\R{円}{えん}\R{玉}{だま} 4まい、 1\R{円}{えん}\R{玉}{だま} 3まい。 なん\R{円}{えん}？\par\vspace{3pt}
\begin{bkfigure}
\begin{tikzpicture}[x=1mm, y=1mm]
  \bkbill{(8,0)}{1000}
  \foreach \x in {21,29.5} {\bkcoin{(\x,0)}{100}}
  \foreach \x in {39,47.5} {\bkcoin{(\x,3.8)}{10}\bkcoin{(\x,-3.8)}{10}}
  \foreach \x in {56,63} {\bkcoin{(\x,3.6)}{1}}
  \bkcoin{(59.5,-3.6)}{1}
\end{tikzpicture}
\end{bkfigure}
\kangaekata 1000が 1つ、 100が 2つ、 10が 4つ、 1が 3つ。 \R{左}{ひだり}から \R{千}{せん}・\R{百}{ひゃく}・\R{十}{じゅう}・\R{一}{いち}の くらいに かく。

\kotae 1243\R{円}{えん}\par\vspace{3pt}
\noindent ぎゃくに、 1243を くらいに わけると、 1000が 1つ、 100が 2つ、 10が 4つ、 1が 3つ。 \R{左}{ひだり}から じゅんに、 1つずつ よめば よい。
\end{reidai}

\begin{jibun}
1000\R{円}{えん}さつ 2まい、 100\R{円}{えん}\R{玉}{だま} 1まい、 10\R{円}{えん}\R{玉}{だま} 3まい、 1\R{円}{えん}\R{玉}{だま} 6まい。 なん\R{円}{えん}？\par\vspace{3pt}
\begin{bkfigure}
\begin{tikzpicture}[x=1mm, y=1mm]
  \foreach \x in {8,26} {\bkbill{(\x,0)}{1000}}
  \bkcoin{(39,0)}{100}
  \foreach \x in {48,56.5,65} {\bkcoin{(\x,0)}{10}}
  \foreach \x in {74,81,88} {\bkcoin{(\x,3.6)}{1}\bkcoin{(\x,-3.6)}{1}}
\end{tikzpicture}
\end{bkfigure}
\bkkotaeran[6\zw]
\bkjans{2136\R{円}{えん}}
\end{jibun}

\begin{jibun}
2507は、 1000が\enspace\bkkakikomi\enspace つ、 100が\enspace\bkkakikomi\enspace つ、\par
\noindent 10が\enspace\bkkakikomi\enspace つ、 1が\enspace\bkkakikomi\enspace つ。
\bkjans{1000が 2つ、 100が 5つ、 10が 0（ひとつも ない）、 1が 7つ}
\end{jibun}

\begin{column}{どうして 10ずつ なの？}
\noindent\begin{minipage}[c]{0.68\linewidth}
りょうての ゆびが、 5\R{本}{ほん}と 5\R{本}{ほん}で 10\R{本}{ぽん} だからと いわれて います。 とけいの ように、 60ずつ まとめる かぞえかたも あります。
\end{minipage}\hfill
\begin{minipage}[c]{0.3\linewidth}\centering
\begin{tikzpicture}[x=1.3mm, y=1.3mm]
  \bkqbhand{(-8,0)}{-1}\bkqbhand{(8,0)}{1}
  \node[font=\sffamily\bfseries\small, text=bkTheme] at (-8,-8) {5};
  \node[font=\sffamily\bfseries\small, text=bkTheme] at (8,-8) {5};
\end{tikzpicture}
\end{minipage}
\end{column}





\clearpage
%% ---------------- p.32 まちがえやすいところ ----------------
\bksection{tsumazuki}{まちがえやすい ところ}

\begin{tsumazuki}{「にじゅうさん」の かきかた}
\bkbatsu\enspace 203\hspace{3\zw}\bkmaru\enspace 23\par
203は 「にひゃくさん」。 「にじゅうさん」は、 \R{十}{じゅう}の くらいに 2、 \R{一}{いち}の くらいに 3を かきます。
\end{tsumazuki}

\begin{tsumazuki}{くらいの いれちがい}
10が 3つと 1が 5つ。\hspace{1\zw}\bkbatsu\enspace 53\hspace{3\zw}\bkmaru\enspace 35\par
10が いくつ あるかは、 \R{左}{ひだり}の \R{十}{じゅう}の くらいに かきます。
\end{tsumazuki}

\begin{tsumazuki}{\R{千}{せん}の くらいが ある かず}
「せん にひゃく さんじゅう」を かくと……\par
\bkbatsu\enspace 100020030\hspace{3\zw}\bkmaru\enspace 1230\par
きこえた じゅんに 1000、 200、 30と つなげて かかない。 1000が 1つ、 100が 2つ、 10が 3つ、 ばらは なし。 くらいの へやに 1つずつ かいて 1230です。
\end{tsumazuki}

\begin{tsumazuki}{39の つぎ}
\bkbatsu\enspace さんじゅうじゅう\hspace{3\zw}\bkmaru\enspace 40（よんじゅう）\par
ばらが 10こに なったら たばねて、 10が 4つ。
\end{tsumazuki}



\clearpage
%% ---------------- p.33 つかわれているところ ----------------
\begin{butsuri}
くらいの しくみは、 きかいでも はたらいて います。
\begin{bkfigure}
\begin{tikzpicture}[x=1mm, y=1mm]
  % くるまの メーター
  \fill[white, rounded corners=3mm] (0,0) rectangle (59,46);
  \node[font=\sffamily\bfseries\small] at (29.5,41) {くるまの メーター};
  \bkqbodo{(10.5,23)}{1,2,3,9}{0}
  \bkqbodo{(10.5,3)}{1,2,4,0}{3}
  \draw[bk arrow, bkTheme] (29.5,21.5) -- (29.5,17.5);
  \node[anchor=west, font=\sffamily\bfseries\fontsize{8.5}{9}\selectfont, text=bkTheme] at (31,19.5) {1 すすむと};
  % 電卓
  \begin{scope}[shift={(63,0)}]
  \fill[white, rounded corners=3mm] (0,0) rectangle (59,46);
  \node[font=\sffamily\bfseries\small] at (29.5,41) {でんたく};
  \bkdisplay{(6,23)}{22}{2}
  \bkdisplay{(32,23)}{22}{23}
  \draw[bk arrow, bkTheme] (28.5,28.5) -- (31.5,28.5);
  \foreach \t/\c/\n in {1/bkPaper/,2/bkOrange!35/①,3/bkOrange!35/②} {
    \node[draw=bkGray, fill=\c, rounded corners=0.8mm, minimum width=8mm, minimum height=6.5mm, inner sep=0pt, font=\sffamily\bfseries\small] at (8.5+\t*10.5,13) {\t};
    \node[font=\sffamily\bfseries\fontsize{8.5}{9}\selectfont, text=bkOrange!85!black] at (8.5+\t*10.5,5.5) {\n};}
  \end{scope}
\end{tikzpicture}
\end{bkfigure}
\begin{itemize}[label={\color{bkPurple}●}, leftmargin=1.4\zw, itemsep=0pt, topsep=0pt]
\item メーターの \R{右}{みぎ}の わは、 9から 0に もどる とき、 となりの わを 1つ すすめます。
\item でんたくで 2、 3と おすと、 2は \R{左}{ひだり}へ うごいて \R{十}{じゅう}の くらいに なります。
\item コンピュータは、 でんきが 「ながれる」「ながれない」の 2とおりで かずを あらわすので、 2ずつ たばねて います。
\end{itemize}
\end{butsuri}

\bkfukidashi{delta}{10こで 1つに まとめて、 それを また 10こで まとめる。 くりかえすと、 どんどん \R{大}{おお}きな かずに なるよ。 \R{下}{した}の \R{路線}{ろ|せん}の はじまりだよ！}

\begin{bkfigure}
\begin{tikzpicture}[x=1mm, y=1mm]
  \draw[bkRoute3, line width=3pt] (10,0) -- (128,0);
  \node[anchor=west, inner sep=0pt] at (0,12) {\bkrouteicon[1.4em]{3}\hspace{0.3em}{\sffamily\bfseries\small\color{bkRoute3!70!black}\R{倍々}{ばい|ばい}の \R{路線}{ろ|せん}}};
  \foreach \x/\c/\t/\s in {10/bkVol1/{1-2}/{10ずつ},49.3/bkVol2/{2-7}/{るいじょう},88.7/bkVol3/{3-11}/{しすう},128/bkVol5/{5-14}/{e}} {
    \fill[white, draw=\c, line width=1.6pt] (\x,0) circle (2.6);
    \node[anchor=base, font=\sffamily\bfseries\footnotesize, text=\c!85!black] at (\x,-8.4) {\t};
    \node[anchor=base, font=\sffamily\footnotesize, text=bkInk] at (\x,-15.8) {\s};}
  \node[fill=bkTheme, text=white, rounded corners=1mm, font=\sffamily\bfseries\fontsize{8}{9}\selectfont, inner sep=1.5pt] at (10,5.5) {いま ここ};
\end{tikzpicture}
\end{bkfigure}

\clearpage
%% ---------------- p.34〜35 もんだい ----------------
\bkendproblems

\bkproblemgroup{きほん}{1}

\bkproblem なん\R{本}{ぼん} あるかな？
\par\nobreak\vspace{2pt}
\begin{bkfigure}
\begin{tikzpicture}[x=1mm, y=1mm]
  \foreach \px/\l in {0/（1）,64/（2）} {\fill[white, draw=bkLine, line width=0.8pt, rounded corners=2.5mm] (\px,0) rectangle ++(60,27);
    \node[anchor=north west, font=\small, text=bkInk] at (\px+0.5,26.5) {\l};
    \node[anchor=north, inner sep=0pt] at (\px+30,-2) {\bkkakikomi\enspace{\small \R{本}{ほん}}};}
  % （1）2たばと 7本
  \bkqbbundle{(6.8,2.5)}{16}{bkRoute4}\bkqbbundle{(22.3,2.5)}{16}{bkRoute4}
  \foreach \k in {0,...,6} {\bkqbstraw{(38.3+\k*2.5,10.5)}{16}{90}{bkRoute4}}
  % （2）3たばと 4本
  \foreach \x in {68.2,82.7,97.2} {\bkqbbundle{(\x,2.5)}{16}{bkRoute2}}
  \foreach \k in {0,...,3} {\bkqbstraw{(112.7+\k*2.4,10.5)}{16}{90}{bkRoute2}}
\end{tikzpicture}
\end{bkfigure}
\bkpans{（1）27\R{本}{ほん}\quad（2）34\R{本}{ほん}}

\bkproblem あいて いる ますに かずを かこう。\par\vspace{2pt}
\noindent\hspace*{1\zw}（1）10が 6つと 1が 8つで\enspace\bkkakikomi\par\vspace{4pt}
\noindent\hspace*{1\zw}（2）75は、 10が\enspace\bkkakikomi\enspace つと 1が\enspace\bkkakikomi\enspace つ
\bkpans{（1）68\quad（2）10が 7つと 1が 5つ}

\bkproblem あいて いる ますに はいる かずを かこう。\par\vspace{2pt}
\noindent\hspace*{1\zw}（1）28、\enspace 29、\enspace\bkkakikomi、\enspace 31\hspace{3\zw}（2）57、\enspace 58、\enspace 59、\enspace\bkkakikomi
\bkpans{（1）30\quad（2）60}

\bkproblem なん\R{円}{えん}かな？
\par\nobreak\vspace{2pt}
\noindent\begin{minipage}[c]{0.6\linewidth}\centering
\begin{tikzpicture}[x=1.25mm, y=1.25mm]
  \foreach \x in {0,8.5,17} {\bkcoin{(\x,0)}{100}}
  \bkcoin{(28,0)}{10}
  \foreach \x in {38,45,52} {\bkcoin{(\x,3.6)}{1}\bkcoin{(\x,-3.6)}{1}}
\end{tikzpicture}
\end{minipage}\hfill
\begin{minipage}[c]{0.38\linewidth}
\bkkotaeran[6\zw]
\end{minipage}
\bkpans{316\R{円}{えん}}

\clearpage
\bkproblem あいて いる ますに かずを かこう。\par\vspace{2pt}
\noindent\hspace*{1\zw}（1）1000が 2つ、 100が 5つ、 10が 1つ、 1が 8つで\enspace\bkkakikomi[15mm]\par\vspace{4pt}
\noindent\hspace*{1\zw}（2）3462は、 1000が\enspace\bkkakikomi\enspace つ、 100が\enspace\bkkakikomi\enspace つ、\par
\noindent\hspace*{3\zw}10が\enspace\bkkakikomi\enspace つ、 1が\enspace\bkkakikomi\enspace つ
\bkpans{（1）2518 （2）1000が 3つ、 100が 4つ、 10が 6つ、 1が 2つ}

\bkproblemgroup{チャレンジ}{1}

\bkproblem あいて いる ますに はいる かずを かこう。\par\vspace{2pt}
\noindent\hspace*{1\zw}（1）97、\enspace 98、\enspace 99、\enspace\bkkakikomi\hspace{3\zw}（2）118、\enspace 119、\enspace\bkkakikomi
\bkpans{（1）100\quad（2）120}

\bkproblem 1000\R{円}{えん}さつ 2まい、 100\R{円}{えん}\R{玉}{だま} 3まい、 10\R{円}{えん}\R{玉}{だま} 1まい、 1\R{円}{えん}\R{玉}{だま} 5まい。 なん\R{円}{えん}？
\par\nobreak\vspace{2pt}
\begin{bkfigure}
\begin{tikzpicture}[x=1.25mm, y=1.25mm]
  \foreach \x in {8,26} {\bkbill{(\x,0)}{1000}}
  \foreach \x in {40.5,48.5,56.5} {\bkcoin{(\x,0)}{100}}
  \bkcoin{(66,0)}{10}
  \foreach \x in {74.5,81} {\bkcoin{(\x,3.6)}{1}\bkcoin{(\x,-3.6)}{1}}
  \bkcoin{(87.5,0)}{1}
\end{tikzpicture}
\end{bkfigure}
\bkkotaeran[6\zw]
\bkpans{2315\R{円}{えん}}

\begin{furikaeri}
\begin{checklist}
\item 10ずつ たばねて かぞえられる。\hfill{\footnotesize\sffamily\color{bkTheme}→\pageref{bk:1-2:tabaneru}ページ}
\item 23の 2と 3の いみが わかる。\hfill{\footnotesize\sffamily\color{bkTheme}→\pageref{bk:1-2:kurai}ページ}
\item 100、 1000、 10000が わかる。\hfill{\footnotesize\sffamily\color{bkTheme}→\pageref{bk:1-2:hyaku}ページ}
\end{checklist}
チェックが つかなかったら、 →の ページに もどろう。
\end{furikaeri}

\clearpage
